\section{Source register for admitted statements}
\label{sec:source-register}

\begingroup
\sloppy
\hbadness=2000

This working register lists only sources already used in mathematical text.
Machine-readable identities, hashes, manifestations, and reading status are
maintained in \path{state/source_registry.jsonl}.
\begin{thebibliography}{99}

\bibitem{Crandall1978}
R.~E. Crandall,
\emph{On the ``$3x+1$'' problem},
Mathematics of Computation \textbf{32} (1978), no.~144, 1281--1292.  All
twelve printed pages were read from the published PDF.  The full theorem,
proof, defect, repair, and nonclaim record is maintained in
\path{qa/CRANDALL-1978-F010-full-paper-audit.md}; deterministic finite
arithmetic checks are in \path{certificates/crandall_1978_checks.py}.
Crandall's Theorem~7.2 is printed for $n\geq4$; its proof explicitly assumes
$k\leq q_n$ but invokes a lemma with strict denominator range, so the edition
  supplies the missing $k=q_n$ justification.  Printed p.~1282 routes the
  historical verification below $10^9$ to references~1 and~5, but content
  and metadata inspection shows that both refer to Conway's same 1972 paper.
  Crandall's reference~1 misprints 49--52 as 45--52; reference~5 misprints
  \emph{Zbl}~0337.10041 as ``Band~233, p.~10041,'' although $10041$ is an
  accession suffix rather than a page.  Printed p.~1290 uses the resulting
  self-return exclusion without a new inline citation.  The published report
  is now content-read; its underlying computation remains unlocated and the
  numerical corollary remains conditional on that report.

\bibitem{Conway1972}
J.~H. Conway,
\emph{Unpredictable iterations},
in \emph{Proceedings of the 1972 Number Theory Conference}, University of
Colorado, Boulder, 1972, 49--52; \emph{Zbl}~0337.10041.  Reprinted in
Jeffrey C. Lagarias, editor, \emph{The Ultimate Challenge: The $3x+1$
Problem}, American Mathematical Society, 2010, 219--223.  Conway's original
text occupies reprint pp.~219--221; editorial commentary begins on p.~221.
All five reprint pages were inspected at 300 dpi, with the copyrighted
third-party extract retained only as a local inspection manifestation and
not represented as an official open-access or redistribution-authorized
copy.  On p.~219 Conway defines the ordinary one-step map and reports
verification for $n\leq10^9$ by D.~H. Lehmer, Emma Lehmer, and
J.~L. Selfridge.  No underlying computational record or reproducibility data
are supplied.  Exact identity, manifestation, clock transfer, and nonclaim
boundaries are recorded in
\path{qa/CONWAY-1972-F022-historical-verification-audit.md}.

\bibitem{AbramowitzStegun1964}
Milton Abramowitz and Irene A. Stegun, editors,
\emph{Handbook of Mathematical Functions with Formulas, Graphs, and
Mathematical Tables},
National Bureau of Standards Applied Mathematics Series~55, U.S. Government
Printing Office, Washington, D.C., issued June 1964.  Crandall cites the
ninth Dover printing, New York, 1965; the complete controlling local witness
is instead the NBS first-printing facsimile and is not represented as the
cited manifestation.  The relevant definitions and formulas are
6.5.1--6.5.3, 6.5.11, 6.5.13, 6.5.34--6.5.35, 26.2.1--26.2.6, 26.2.29,
26.4.1--26.4.2, 26.4.11, 26.4.19, and 26.4.21; the exact manifestation,
locator, morphism, and nonclaim boundaries are recorded in
\path{qa/ABRAMOWITZ-STEGUN-1964-F013-poisson-audit.md}.

\bibitem{Pillai1931}
S.~Sivasankaranarayana Pillai,
\emph{On the inequality $0<a^x-b^y\leq n$},
Journal of the Indian Mathematical Society \textbf{19} (1931), 1--11.
All eleven printed pages were read from the controlling article scan.  Its
title, author, journal pagination, complete page sequence, and exact
bibliographic entries agree, but the witness is not represented as a
publisher-hosted download.  The source theorem, false degree assertion,
constant-type defect, residue-class gap, separately attributed repair, and
Crandall crosswalk are recorded in
\path{qa/PILLAI-1931-F014-fixed-difference-audit.md}.

\bibitem{Herschfeld1936}
Aaron Herschfeld,
\emph{The equation $2^x-3^y=d$},
Bulletin of the American Mathematical Society \textbf{42} (1936), no.~4,
231--234, doi:10.1090/S0002-9904-1936-06275-0.  All four printed pages were
read from the official AMS article PDF.  The positive-exponent theorem,
two-sided $|d|$ quantifier, ineffective threshold, two order coordinates,
Pillai dependency repair, separately attributed all-integer extension,
source label defect, unpublished-computation boundary, and distinct
Siegel-1929 dependency are recorded in
\path{qa/HERSCHFELD-1936-F018-eventual-uniqueness-audit.md}; deterministic
finite checks are in \path{certificates/herschfeld_1936_checks.py}.  Crandall's
reference~10 spelling ``Herschefeld'' is not the spelling on the primary
byline.  The conditional primitive-root paragraph printed on p.~1288 is
reconstructed from the full monotone return equation, with its exact domain
and two residue branches, in
\path{qa/CRANDALL-1978-F020-primitive-root-audit.md}; the finite selector
certificate is \path{certificates/primitive_root_return_checks.py}.

\bibitem{Polya1918}
Georg P\'olya,
\emph{Zur arithmetischen Untersuchung der Polynome},
Mathematische Zeitschrift \textbf{1} (1918), no.~2--3, 143--148,
doi:10.1007/BF01203608.  Every article page was read from the official SUB
G\"ottingen/GDZ scan.  Satz~I, its explicit Thue dependency, and the exact
fixed-prime gap theorem on printed pp.~147--148 are audited in
\path{qa/POLYA-1918-F015-smooth-gap-audit.md}.

\bibitem{Thue1908I}
Axel Thue,
\emph{Bemerkungen \"uber gewisse N\"aherungsbr\"uche algebraischer Zahlen},
Videnskabs-Selskabets Skrifter, I. Math.-Naturv. Klasse, 1908, no.~3,
34 pp.  The complete 1908 academy volume was read from the public-domain
BHL/Internet Archive scan.  Satz~12 is on printed pp.~30--31 / physical PDF
pp.~182--183.

\bibitem{Thue1908II}
Axel Thue,
\emph{Om en generel i store hele tal ul{\o}sbar ligning},
Videnskabs-Selskabets Skrifter, I. Math.-Naturv. Klasse, 1908, no.~7,
1--15.  The opening theorem is on printed p.~1 / physical PDF p.~335 of the
same academy-volume scan.  Both located primary scans omit printed pp.~2--3;
therefore P\'olya's locator to p.~3 is retained as a visible manifestation
gap, not represented as directly read.

\bibitem{Thue1909}
Axel Thue,
\emph{\"Uber Ann\"aherungswerte algebraischer Zahlen},
Journal f\"ur die reine und angewandte Mathematik \textbf{135} (1909),
284--305, doi:10.1515/crll.1909.135.284.  Theorem~IV was read on printed
pp.~303--304 / PDF pp.~21--22 of the official SUB G\"ottingen/GDZ scan.  The
three identities, locators, exact exceptional-form dependency, quadrant map,
residue-cell cubic, and scan gap are audited in
\path{qa/THUE-1908-1909-F017-binary-form-audit.md}.

\bibitem{Siegel1921}
C.~L. Siegel,
\emph{Approximation algebraischer Zahlen},
Mathematische Zeitschrift \textbf{10} (1921), no.~3--4, 173--213,
doi:10.1007/BF01211608.  The complete article and the controlling theorem
pages were read from the official GDZ scan.  The introductory denominator
form, Hauptsatz, proof route, Zus\"atze, height-to-denominator map, exact
constant type, and actual-degree use are recorded in
\path{qa/SIEGEL-1921-F016-approximation-audit.md}.

\bibitem{Everett1977}
C.~J. Everett,
\emph{Iteration of the number-theoretic function
$f(2n)=n$, $f(2n+1)=3n+2$},
Advances in Mathematics \textbf{25} (1977), no.~1, 42--45,
doi:10.1016/0001-8708(77)90087-1.  The publisher identity was verified by
DOI and PII.  All four printed pages were read from a matching VOR facsimile;
the publisher-hosted PDF hash remains unavailable.  The complete theorem,
proof, manifestation, crosswalk, and nonclaim audit is in
\path{qa/EVERETT-1977-F012-density-audit.md}; deterministic finite checks are
in \path{certificates/everett_1977_checks.py}.

\bibitem{Steuding2005}
J\"orn Steuding,
\emph{Diophantine Analysis},
Discrete Mathematics and Its Applications, CRC Press/Taylor \& Francis
Group, 2005; eBook PDF version dated 13 July 2011,
ISBN 978-1-4200-5720-1.  The continued-fraction definitions, recurrence and
determinant identities, exact convergent-error formula, law of best
approximations, unimodular equivalence, Serret theorem, Markoff formula, and
starred Exercise~5.22 were read from the local published-book PDF.  The
Exercise~5.22 solution, printed matrix defects, exact clock crosswalk, and
nonclaims are recorded in
\path{qa/STEUDING-2005-F011-exercise-522-audit.md}; exact finite
quadratic-field checks are in
\path{certificates/steuding_522_checks.py}.  The exact transfer from
Theorems~3.1, 3.2, 3.6, and 3.8 to Crandall's Section~7 inequalities, including
the strict $0<y<q_n$ range and the $y=q_n$ equality boundary, is recorded in
\path{qa/STEUDING-2005-F021-crandall-continued-fraction-audit.md}; its finite
arithmetic and pinned-source checks are in
\path{certificates/continued_fraction_dependency_checks.py}.  This independent
support is not an identification of Crandall's cited Roberts or Khinchin
manifestations.

\bibitem{BohmSontacchi1978}
Corrado B\"ohm and Giovanna Sontacchi,
\emph{On the existence of cycles of given length in integer sequences like
$x_{n+1}=x_n/2$ if $x_n$ is even, and $x_{n+1}=3x_n+1$ otherwise},
Atti della Accademia Nazionale dei Lincei. Classe di Scienze Fisiche,
Matematiche e Naturali. Rendiconti, series~8, \textbf{64} (1978), 260--264.

\bibitem{BernsteinLagarias1996}
Daniel J. Bernstein and Jeffrey C. Lagarias,
\emph{The $3x+1$ conjugacy map},
Canadian Journal of Mathematics \textbf{48} (1996), no.~6, 1154--1169,
doi:10.4153/CJM-1996-060-x.

\bibitem{Bernstein1994}
Daniel J. Bernstein,
\emph{A non-iterative $2$-adic statement of the $3N+1$ conjecture},
Proceedings of the American Mathematical Society \textbf{121} (1994),
no.~2, 405--408, doi:10.1090/S0002-9939-1994-1186982-9.  The official AMS
PDF and the author-hosted PDF were read in full; the displayed formulas were
checked visually because the AMS text extraction corrupts symbols.

\bibitem{Lagarias1985}
Jeffrey C. Lagarias,
\emph{The $3x+1$ problem and its generalizations},
American Mathematical Monthly \textbf{92} (1985), no.~1, 3--23,
doi:10.1080/00029890.1985.11971528.  The complete JSTOR page-image scan was
read; the author-linked 1996 web conversion is an auxiliary reprint rather
than the controlling locator witness.

\bibitem{Lagarias1990}
Jeffrey C. Lagarias,
\emph{The set of rational cycles for the $3x+1$ problem},
Acta Arithmetica \textbf{56} (1990), no.~1, 33--53,
doi:10.4064/aa-56-1-33-53.  The complete official IMPAN scan was read.

\bibitem{Stanley2024}
Richard P. Stanley,
\emph{Enumerative Combinatorics}, volume~2, second edition,
Cambridge Studies in Advanced Mathematics, volume~208,
Cambridge University Press, 2024,
doi:10.1017/9781009262538.  Exercise~7.89(a), equations~(7.149)--(7.150)
on printed p.~480, and the solution on printed p.~555 were inspected for the
fixed-content primitive-necklace formula.

\bibitem{LaarhovenDeWeger2012}
Thijs Laarhoven and Benne de Weger,
\emph{The Collatz conjecture and De Bruijn graphs},
arXiv:1209.3495v1 (2012).  The complete source TeX was inspected.  The
official v1 PDF's visible 27 November 2024 date comes from
\verb|\date{\today}| and is not a mathematical revision.  Exact audit details
are in the machine-readable source record.

\bibitem{Urata2003}
Toshio Urata,
\emph{The Collatz problem over $2$-adic integers},
The Bulletin of Aichi University of Education \textbf{52} (Natural Science)
(2003), 5--11.  The official Aichi repository PDF was inspected in full; two
later local re-typesettings are retained only as unauthenticated editorial
derivatives and do not control attribution.

\bibitem{Matthews2010AuthorForm}
K.~R. Matthews,
\emph{Generalized $3x+1$ mappings: Markov chains and ergodic theory},
in Jeffrey C. Lagarias (ed.), \emph{The Ultimate Challenge: The $3x+1$
Problem}, American Mathematical Society, 2010, 79--103.  The inspected local
author form explicitly contains small additions and deletions and is not
treated as text-identical to the published chapter.

\bibitem{MatthewsWatts1984}
K.~R. Matthews and A.~M. Watts,
\emph{A generalization of Hasse's generalization of the Syracuse algorithm},
Acta Arithmetica \textbf{43} (1984), no.~2, 167--175,
doi:10.4064/aa-43-2-167-175.  Official IMPAN scan inspected.

\bibitem{KrasikovLagarias2003}
Ilia Krasikov and Jeffrey C. Lagarias,
\emph{Bounds for the $3x+1$ problem using difference inequalities},
Acta Arithmetica \textbf{109} (2003), no.~3, 237--258,
doi:10.4064/aa109-3-4.  The complete 22-page official IMPAN scan and the
complete arXiv source \texttt{math/0205002v1} were read.  The published
Theorem~6.1 gives
\(\pi_a(x)\ge x^{0.84}\) beyond an \(a\)-dependent threshold for every
positive \(a\not\equiv0\pmod3\).  Its proof reports a computer-found feasible
\(k=11\) point at \(\lambda=1.7922310\), but the paper supplies neither that
coordinate vector nor code or an exact feasibility certificate.  The
published theorem is retained; independent reproducibility of this finite
premise remains open.  The exact source objects and proof chain, six
published source defects, two preprint-only defects corrected in the
journal, cycle-target handoff, shortened/unshortened clock morphism, and
nonclaims are recorded in
\texttt{qa/KRASIKOV-LAGARIAS-2003-F029-difference-inequality-audit.md} and
Subsection~\ref{subsec:KL-predecessor-clock}.  Pinned-source, exact-exponent,
and finite clock checks are in
\texttt{certificates/krasikov\_lagarias\_2003\_checks.py}.

\bibitem{Terras1976}
Riho Terras,
\emph{A stopping time problem on the positive integers},
Acta Arithmetica \textbf{30} (1976), no.~3, 241--252,
doi:10.4064/aa-30-3-241-252.  Every printed page was read from the official
IMPAN scan.  The source uses the shortened map on
\(\mathbb N=\{0,1,2,\ldots\}\), defines the chronological parity vector,
coefficient stopping time \(\tau\), and actual stopping time \(\chi\), and
proves the exact residue-fibre theorem and coefficient-tail estimate.  Four
local defects are retained rather than normalized: Proposition~1.7 omits the
hypothesis \(\lambda_k(n)<1\); Theorem~1.9 cites Proposition~1.9 rather than
1.8; Definition~1.12 prints the proper-prefix range through \(k-2\) rather
than \(k-1\); and Theorem~1.17 drops \(+1\) from its terminal-string bound.
The exact counterexamples, repairs, and finite checks are recorded in
\texttt{qa/TERRAS-ALLOUCHE-KOREC-F030-density-lineage-audit.md} and
\texttt{certificates/terras\_allouche\_korec\_checks.py}.  The 1976 text's
point-mass and coefficient-tail results do not by themselves replace the
finite-complement density-existence proof supplied by Terras's 1979 note.

\bibitem{Terras1979}
Riho Terras,
\emph{On the existence of a density},
Acta Arithmetica \textbf{35} (1979), no.~1, 101--102,
doi:10.4064/aa-35-1-101-102.  Both printed pages were read from the official
IMPAN scan.  The note records doubts raised by A.~Garsia, H.~Moeller, and the
editors, then proves for every fixed positive \(k\) the exact identity
\(\delta[\chi\geq k]=\mathbb P[\tau\geq k]\) by finite additivity and
complementation.  This closes the density-existence dependency left implicit
by ``forming sums'' in the 1976 paper; no infinite sum of densities is used.

\bibitem{Allouche1979}
Jean-Paul Allouche,
\emph{Sur la conjecture de ``Syracuse--Kakutani--Collatz''},
S\'eminaire de Th\'eorie des Nombres de Bordeaux, ann\'ee 1978--1979,
expos\'e no.~9, printed pp.~9-01--9-15, CNRS, Talence, 1979,
MR0567874.  All fifteen printed pages were read from the GDZ manifestation
identified by \texttt{PPN320141322\_0008/LOG\_0014}.  The edition preserves
the source's unnormalized data \(n,d,A_d,\varphi,g\), the exact affine fibre
coordinate, bounded remainder, binomial branch counts, and quantitative
distribution theorem.  The strict exponent
\(3/2-\log_3 2\) is reconstructed from those lemmas and is not presented as
a decimal theorem printed by Allouche.  The displayed definition of the
analogue map \(F:\mathbb Z\to\mathbb Z\) in Theorem~2 omits some divisible
residue classes (for example \(6\bmod8\) when \(d=2,u=3\)); the live reader
states the unique bounded residue correction as a separately attributed
repair before using the same-slopes/opposite-dynamics obstruction.

\bibitem{Korec1994}
Ivan Korec,
\emph{A density estimate for the $3x+1$ problem},
Mathematica Slovaca \textbf{44} (1994), no.~1, 85--89,
MR1290275, Zbl~0797.11027.  Every printed page was read from the DML-CZ
manifestation at \texttt{hdl.handle.net/10338.dmlcz/133225}.  Theorem~1 proves
natural density one for
\(\{y:\exists n,\ T^n(y)<y^c\}\) at every strict
\(c>\log_4 3\), with witness \(n\leq\log_2 y\) on the retained range.  Its
proof contains three locally repairable defects: the auxiliary parameter
prints \(\log_4 3\) where the required denominator is \(\log_2 3\); the
claimed logarithmic equivalence omits a factor \(\log_2 3\) in one
denominator (the printed condition remains stronger and sufficient); and the
proof text uses a strict parity-count boundary while \(U(m,d)\) and the final
block count use a non-strict one.  The theorem is used only after all three
repairs are displayed.  No endpoint exponent, uniform iterate, explicit
density rate, all-integer statement, or Collatz conclusion is attributed to
the source.

\bibitem{Kohl2008}
Stefan Kohl,
\emph{Algorithms for a class of infinite permutation groups},
Journal of Symbolic Computation \textbf{43} (2008), no.~8, 545--581,
doi:10.1016/j.jsc.2007.12.001.  The definition, fixed coefficient
representation, and image-density remark were inspected in the local author
PDF; the later group algorithms have not yet been certified in this edition.

\bibitem{Siegel2007v1}
Maxwell C. Siegel,
\emph{Syracuse Random Variables and the Periodic Points of Collatz-type
maps}, arXiv:2007.15936v1 (historical source dated 13 July 2020).  The
version-pinned source TeX controls only claims attributed to v1 here.  The
definition and one-position series for the numen were read at source
lines~405--568, and the explicit discussion of Tao's finite Syracuse random
variables was read at lines~3441--3501.  This historical version is not
silently identified with the renamed v3 work.

\bibitem{Siegel2007v3}
Maxwell C. Siegel,
\emph{The Collatz Conjecture \& Non-Archimedean Spectral Theory---Part I:
Arithmetic Dynamical Systems and Non-Archimedean Value Distribution Theory},
arXiv:2007.15936v3 (source dated 27 March 2023; revised 24 April 2023).
The finite numen, explicit series, and Syracuse-variable passages were read
at source lines~1524--1582, 1851--2000, and 3083--3184.  The exact
run-length/Haar/Syracuse crosswalk is proved in
Proposition~\ref{prop:Siegel-Tao-Syracuse-coordinate}; the three quarantined
Fourier-presentation defects are recorded in
Source Issue~\ref{sourceissue:Siegel-2007-v3-Fourier}.  This is a distinct
versioned source from both v1 above and arXiv:2601.17030v3 below.

\bibitem{Siegel2026v3}
Maxwell C. Siegel,
\emph{The Hydra Map and Numen Formalisms for Collatz-Type Problems},
arXiv:2601.17030v3 (2026).  Version-pinned source TeX controls this working
edition; the Hydra, numen, Correspondence Principle, and v3 Fourier boundary
were inspected against the 30-page source PDF.

\bibitem{Tao2026v7}
Terence Tao,
\emph{Almost all orbits of the Collatz map attain almost bounded values},
arXiv:1909.03562v7 (revision 16 July 2026); published in Forum of Mathematics,
Pi \textbf{10} (2022), e12, doi:10.1017/fmp.2022.8.  The version-pinned v7
source TeX controls this working edition; the frozen discovery routes point
to v5 and are not silently identified with either v7 or the journal text.
All v7 mathematical source through its proof ending at line~1859 has been
read at content level; line~1860 is blank, and lines~1863--1952 are the
active bibliography.  The invalid additive
endpoint buffer at lines 758--763 is replaced by the proved multiplicative
interval in Proposition~\ref{prop:Tao-v7-endpoint-buffer}; its dependency use
has been checked through line 925, with the signed lower bound at
lines 790--793 and the malformed line 889 estimate also repaired.  The event
chain, fixed-\(M\) affine equivalence, positive valuation domain, the distinct
line 868 and line 899 previous-iterate coefficients, line 917 probability
transfer, the line 689 \(1.9n_0\) index, and the non-uniform line 913
summation are repaired in
Proposition~\ref{prop:Tao-v7-Prop52-cn-repair}.  This certifies source
Proposition 5.2 and the uniform \(c_n\) lemma, and checks their Section 5
dependency use through line 925.  The exact audit is in
\texttt{qa/TAO-V7-F001-proof-text-audit.md}; the pinned-source, locator, and
arithmetic checks are in
\texttt{certificates/tao\_v7\_endpoint\_buffer\_checks.py} and
\texttt{certificates/tao\_v7\_prop52\_cn\_checks.py}.  Fine-scale mixing,
Fourier decay, and renewal are reconstructed in
Propositions~\ref{prop:Tao-v7-Fourier-to-mixing}--
\ref{prop:Tao-v7-renewal-monotonicity} and
Corollary~\ref{cor:Tao-v7-fine-scale-branch}; the line-level and
version-comparison record is
\texttt{qa/TAO-V5-V7-JOURNAL-F024-fourier-renewal-audit.md}; the exact finite
lift, offset, quotient-frequency, holding-law, and entry-budget kernels are in
\texttt{certificates/tao\_v7\_fourier\_renewal\_checks.py}.  This closes the
v7 fine-scale branch through Proposition~1.14 at its printed uniform scope.
The independent residue-fibre proof of Proposition~1.9 is
Proposition~\ref{prop:Tao-v7-Prop19-valuation-law}; its three-manifestation
comparison is
\texttt{qa/TAO-V5-V7-JOURNAL-F025-prop19-valuation-audit.md}, and its exact
finite kernels are checked by
\texttt{certificates/tao\_prop19\_valuation\_checks.py}.  The exact source
edge from Propositions~1.14 and~1.9 to Proposition~1.11 is reconstructed in
Theorem~\ref{thm:Tao-v7-Prop111-first-passage}.  Its three-manifestation
comparison, source-defect boundary, affine path--endpoint bijection, integer
window, and exact mixing collapse are recorded in
\texttt{qa/TAO-V5-V7-JOURNAL-F026-prop111-first-passage-audit.md}; its finite
integration kernels are checked by
\texttt{certificates/tao\_prop111\_transport\_checks.py}.  Proposition~1.11
is therefore closed at its printed v7 scope.  Its exact deductions through
Theorem~3.1 and Theorem~1.6 to Theorem~1.3, including the indexed
nested-threshold morphism, fixed-tail repair, dyadic stratum bijection and
harmonic transport, are reconstructed in
Subsection~\ref{subsec:Tao-main-reduction}.  The three-manifestation defect
record is
\texttt{qa/TAO-V5-V7-JOURNAL-F027-main-theorem-reduction-audit.md}, and the
finite exact kernels are checked by
\texttt{certificates/tao\_main\_reduction\_checks.py}.  The remaining
whole-version comparison is recorded in
\texttt{qa/TAO-V5-V7-JOURNAL-F028-whole-version-citation-audit.md}.
It classifies all 28 non-whitespace v5--v7 hunks, the journal's
15-v7/12-v5/one-hybrid manifestation matrix, the exact terminal boundary,
and all active citations.  Its Section~6 calculation proves that the v5 and
journal half-window is too small for the printed Archimedean argument,
whereas v7's \(0.99\)-window closes; see
Source Issue~\ref{sourceissue:Tao-v5-v7-section-six-reserve}.  The
deterministic source, bibliography, citation, and finite-margin checks are in
\texttt{certificates/tao\_whole\_version\_checks.py}.  The substantive cited
literature not yet present locally and the durable whole-paper release gate
remain open.

\bibitem{Tao2019v5}
Terence Tao,
\emph{Almost all orbits of the Collatz map attain almost bounded values},
arXiv:1909.03562v5 (version-pinned comparison source).  Its source archive
and extracted TeX were admitted to the topical shelf and compared directly
with v7 and the published paper.  Its Proposition~1.11 deduction at source
lines~643--919 has the same substantive modular, logarithmic-mass, affine,
and integer-window calculations, but lines~789--791 interchange passage time
and passage location; v7 lines~795--797 provide the typed correction used by
the edition.  In Section~6, v5's \(\frac12C_A^2\log n\) source window leaves
only \(\frac12\log2\) times that quadratic reserve after exponentiation,
below the optimized cost
\(\log^22/(4\log(4/3))\); it also drops the factor \(\log2\) in the printed
conversion.  An explicit concentration-admissible family makes one claimed
sub-\(3^n\) summand exceed \(3^n\).  This invalidates that printed size
argument without establishing a residue collision or disproving the
standalone conclusion.  V5 also prints
\(1_{\overline E_k\cap B_k\cap C_{k,l}}\) where the factorization requires
\(1_{E_k\cap B_k\cap C_{k,l}}\).  Both defects are repaired in v7 and
reconstructed in
Source Issue~\ref{sourceissue:Tao-v5-v7-section-six-reserve}.  V5 and the
journal also state the stronger
unconditional form of Lemma~7.9, while v7 states an indicator-restricted
form.  The stronger displayed induction has an unresolved \(r=1<R\) gap and
is not used in this edition; see
Source Issue~\ref{sourceissue:Tao-v5-v7-many-triangles}.  After whitespace
deletion, its main reduction at lines~529--587 is exactly equal to v7
lines~535--593; this local equality does not identify the complete versions.
The shared reduction defects and the independently authored exact repair are
recorded in
\texttt{qa/TAO-V5-V7-JOURNAL-F027-main-theorem-reduction-audit.md} and
Subsection~\ref{subsec:Tao-main-reduction}.  The complete source comparison,
including the journal's mixed-version boundary and the exact bibliography,
is F028; local equality of the bibliography and main-reduction blocks still
does not identify the complete manifestations.

\bibitem{Barina2021}
David Barina,
\emph{Convergence verification of the Collatz problem},
The Journal of Supercomputing \textbf{77} (2021), no.~3, 2681--2688,
doi:10.1007/s11227-020-03368-x; published online 1 July 2020.
All eight publisher pages and the distinct eight-page author postprint were
read.  Their SHA-256 hashes are respectively
\nolinkurl{737a225fbc291cc182e940cc7ca31da4d2eb99b64e51a33e380213d8ed6b33a3}
and
\nolinkurl{6df53d8f2a41d032f215209f06a021d35e447f6bdf222f1d557be730b64409dc}.
The map charts, block algorithm, affine sieve coordinate, reported strict
$2^{68}$ bound, record datum, exceptional-start defects, and nonclaims are
reconstructed in Subsection~\ref{subsec:computational-verification} and
\path{qa/BARINA-OLIVEIRA-ROOSENDAAL-F031-computational-verification-audit.md}.

\bibitem{Barina2025}
David Barina,
\emph{Improved verification limit for the convergence of the Collatz
conjecture},
The Journal of Supercomputing \textbf{81} (2025), article~810,
doi:10.1007/s11227-025-07337-0; accepted 21 April 2025 and published online
2 May 2025.  All fourteen publisher pages were read and rendered; the
controlling PDF has SHA-256
\nolinkurl{764fd732ad79545e71440f74bf738bf315b479eb404f68bae2c3b109785a5d06}.
The FIT record's volume~81, no.~1, pp.~1--14 coordinate is retained as a
metadata discrepancy rather than substituted for the publisher's article
number.  Table~10's strict $n<2^{71}$ boundary and 15 January 2025 completion
date control over the surrounding phrase \emph{up to}.  The two reversed
class cardinalities, three exceptional-start cases, optional-kernel and
validation boundaries, and exact four-new/five-cumulative record
reconciliation are stated separately from the reported finite result.

\bibitem{OliveiraSilva2010}
Tomás Oliveira e Silva,
\emph{Empirical Verification of the $3x+1$ and Related Conjectures},
in Jeffrey C. Lagarias, editor,
\emph{The Ultimate Challenge: The $3x+1$ Problem},
American Mathematical Society, 2010, 189--207.
Every printed page was read from a third-party scan whose page images remain
temporary inspection witnesses and are not copied into the topical shelf or
reader.  Identity and pagination were crosschecked against the author's
University of Aveiro bibliography and the AMS volume record.  The chapter's
inclusive $20\cdot2^{58}$, $10^{16}$, and $10^{14}$ reports, duplicate-worker
protocol, generalized maps, finite tables, and expressly non-rigorous Markov
argument are retained at their stated scopes.  Its floor sentence is repaired
locally.  Proposition~2 is not normalized: the exact $i=1,n_0=1$
counterexample and independently proved $6^i$ branch-cylinder replacement
are Source Issue~\ref{sourceissue:Oliveira-proposition-two} and
Proposition~\ref{prop:Oliveira-branch-cylinder}.

\bibitem{Roosendaal2026}
Eric Roosendaal,
\emph{On the $3x+1$ problem},
continuously revised author web page, snapshot identifying itself as last
modified 30 June 2026.  The retained snapshot has SHA-256
\nolinkurl{e3ff21c62b1fe108ecaa9086a2c8679207339b3e086289a159b56099515f157d}.
It was captured over plain HTTP after the HTTPS endpoint presented a
certificate-name failure, so transport authentication is explicitly absent.
The modular-maximum argument is reconstructed in
Proposition~\ref{prop:Roosendaal-mod36}; the page's maximum-variable and
$N<1$ defects are not silently corrected in attribution.

\bibitem{BarinaCode2025}
xbarin02,
\emph{collatz} and \emph{collatz-sieve}, public source-code lineages.
The 2025 article pins \texttt{collatz} commit
\nolinkurl{53c2a0608075d6fe3f10cc6eeeaf50e400c86338}, tree
\nolinkurl{82b9731dcd5108298a2fb2085ba1d76c837de6f5}; the earlier reconstruction
boundary is commit
\nolinkurl{0fabf721434b20f00af820564e3b3781dcabe768}, tree
\nolinkurl{f8597ec3867e5734cd6a9fb57379aef2c670f857}.  Complete history is retained
in a verified Git bundle with SHA-256
\nolinkurl{e7b655972c628f060ecbc7aaffaa381eeebe6868567d14c4f4524c18559cc3dd}.
The separate, article-unpinned sieve dependency is frozen for this audit at
commit \nolinkurl{e5f0c264ba1f5f9c0a7d32a10c3aa61b225c70d1}, tree
\nolinkurl{4b5a3b5afdf35596409c91da7f226e338c7c763a}; its complete bundle has
SHA-256
\nolinkurl{687f00faba696d2a3dc8b0a6f556c9cc6da56756bdb36535ac8f3ed49575983f}.
These pins support code-level reconstruction and the finite
Theorem~\ref{thm:Barina-recursive-k34}; without a completed-work ledger they
do not authenticate a historical distributed run.

\bibitem{BarinaStatus2026}
David Barina's Collatz project,
\emph{verification status API}, dated snapshot generated 28 August 2026 at
19:24:59 $+0200$.  The 198-byte JSON manifestation has SHA-256
\nolinkurl{daf454223878d578dc7a87dc0677aa1925eba602ae295f8547ce25927ce5c5ab}
and reports \texttt{verified\_up\_to\_times\_2\_60=2075},
\texttt{work\_unit\_exp=40}, and first incomplete/unassigned unit
\texttt{2175975546}.  It is a dynamic project assertion, not a paper result
or an independently acquired server ledger.

\bibitem{AizawaIsaacSegar2019}
N. Aizawa, P. S. Isaac, and J. Segar,
\emph{$\mathbb Z_2\times\mathbb Z_2$ generalizations of
$\mathcal N=1$ superconformal Galilei algebras and their representations},
Journal of Mathematical Physics \textbf{60} (2019), 023507,
doi:10.1063/1.5054699; arXiv:1808.09112v1 (submitted 28 August 2018).
The complete 1,127-line source \TeX{} was read.  The paper supplies the color
bracket, its enveloping-algebra construction, complete relation tables,
central extension, decompositions, adjoints, and representations.  It contains
no Collatz map or action; Section~\ref{sec:chatnotes-coefficient-algebra}
uses it only at the exact color-algebra scope and proves the coefficient
crosswalk independently.

\bibitem{DrosteKuichVogler2009}
Manfred Droste, Werner Kuich, and Heiko Vogler, editors,
\emph{Handbook of Weighted Automata},
Springer, 2009, ISBN 978-3-642-01491-8,
eISBN 978-3-642-01492-5.
The content read for Section~\ref{sec:chatnotes-weighted-path-algebra}
comprises physical PDF pages 3--7, 20--22, 27, 84, 89, 134--135, and
191--192: free monoids and morphisms, semirings, formal power series and
polynomials, path-label products, cycle-free and complete-semiring summation
conditions, finite weighted automata, and matrix/formal-series behavior.  The
book supplies this standard framework, not the Collatz-specific theorems of
that section.

\bibitem{KuichSalomaa1986}
Werner Kuich and Arto Salomaa,
\emph{Semirings, Automata, Languages},
EATCS Monographs on Theoretical Computer Science, volume 5,
Springer, 1986.
Physical PDF pages 1--4 of the local chapter-I witness were read for the
definitions of free monoids, monoid morphisms, semirings, formal series,
finite-support polynomials, and Cauchy multiplication.  No Collatz result is
attributed to this source.

\bibitem{Steinberg2016}
Benjamin Steinberg,
\emph{Representation Theory of Finite Monoids},
Universitext, Springer, 2016,
ISBN 978-3-319-43930-3, eBook ISBN 978-3-319-43932-7.
Physical PDF pages 4--5, 22, and 203--204 were read for transformation
monoids, transformation modules, rank, Proposition~13.1, and Lemma~13.2.
The minimum-ideal argument in
Proposition~\ref{prop:finite-residue-minimum-ideal} is proved directly and
then crosswalked against Steinberg's general minimum-rank theorem.

\bibitem{ConnesConsani2016}
Alain Connes and Caterina Consani,
\emph{Geometry of the Arithmetic Site},
\emph{Advances in Mathematics} \textbf{291} (2016), 274--329,
doi:10.1016/j.aim.2015.11.045; arXiv:1502.05580.
The complete local source \TeX{} was used for the present crosswalk.  The
relevant source coordinates are the tensor square and Frobenius maps,
source lines 1151--1214; the semiring
$\operatorname{Conv}(\Z\times\Z)$ and its quotient, lines 1364--1400; and
the reduced-square factorization, lines 1470--1519.  The Frobenius
correspondence construction and its composition law occur at lines
1522--1902, with the three cases stated at lines 1747--1755:
$\Psi(\lambda)\circ\Psi(\lambda')=\Psi(\lambda\lambda')$ when
$\lambda\lambda'$ is irrational; the same equality when both parameters are
positive rational; and
$\Psi(\lambda)\circ\Psi(\lambda')=
\operatorname{Id}_{\varepsilon}\circ\Psi(\lambda\lambda')$ when both
parameters are irrational but their product is positive rational.  In
particular, the
source does not define a Collatz path algebra or identify a nonnegative
bidegree cone with the arithmetic-site square.

\bibitem{BhattLurieWCart}
Bhargav Bhatt and Jacob Lurie,
\emph{The Prismatization of $p$-adic Formal Schemes},
preliminary postscript.
The local source \TeX{} was read at its definition of
$\operatorname{WCart}_X$, lines 480--483; functoriality in formal-scheme
maps, lines 523--530; the separate Witt Frobenius and Hodge--Tate material,
lines 532--550; and the derived/classical warnings, lines 1450--1511.
Section~\ref{sec:symbolic-completion-geometric-boundaries} uses only that
stated functorial scope and independently proves that the raw affine branch
is not an endomorphism of $\operatorname{Spf}(\Z_3)$.

\bibitem{GrosLeStumQuiros2025}
Michel Gros, Bernard Le~Stum, and Adolfo Quir\'os,
\emph{Absolute Calculus and Prismatic Crystals on Cyclotomic Rings},
arXiv:2310.13790v3, 30 June 2025.
The controlling 72-page PDF was read at Theorems~7.5 and 7.12, the de~Rham
and Hodge--Tate comparisons in Sections~8--9, and the prismatic
$F$-crystal/crystalline-representation results in Section~10.  Its base is
the one-variable cyclotomic ring $W[[q-1]]$ with $k$ a perfect field of
characteristic $p>0$ and under its explicit root-of-unity hypotheses.  No two-coordinate Collatz calculus is
attributed to it.

\bibitem{Sarig1999}
Omri M. Sarig,
\emph{Thermodynamic Formalism for Countable Markov Shifts},
\emph{Ergodic Theory and Dynamical Systems} \textbf{19} (1999), no.~6,
1565--1593, doi:10.1017/S0143385799146820.
The official 33-page author preprint was read at physical pages 3, 5--7, 18,
and 24--25.  These pages contain the countable topological Markov shift and
mixing definitions, fixed-base periodic definition of Gurevich pressure,
condition~(13), and Theorem~8.  Condition~(13) is the one-sided big-images
condition: finitely many states $b_1,\ldots,b_n$ receive an edge from every
state.  It is not BIP.  At the theorem's stated topological-mixing, locally
H\"older, and finite-pressure scope, Theorem~8 says that a Gibbs measure
exists exactly when the potential is positive recurrent, its RPF
eigenfunction is uniformly bounded above and below, and condition~(13)
holds; then the Gibbs measure is $h\,d\nu$ and its pressure parameter is the
Gurevich pressure.  This theorem is not silently replaced by the later BIP
criterion.

\bibitem{Sarig2015}
Omri M. Sarig,
\emph{Thermodynamic Formalism for Countable Markov Shifts}, in
\emph{Hyperbolic Dynamics, Fluctuations and Large Deviations},
Proceedings of Symposia in Pure Mathematics \textbf{89}, American
Mathematical Society, 2015, 81--117,
doi:10.1090/pspum/089/01485.
The controlling local file is the 36-page 2015 overview despite its
misleading filename.  Physical pages 14, 16, 18, and 19 were read for
Theorems~4.3, 4.9, 5.3, Definition~5.8, and Theorem~5.9; physical pages
35--36 explicitly list the 1999 article as a distinct source.  The overview
is therefore neither a manifestation nor a theorem locator for the 1999
article.

\bibitem{Sarig2003}
Omri Sarig,
\emph{Existence of Gibbs Measures for Countable Markov Shifts},
Proceedings of the American Mathematical Society \textbf{131} (2003),
1751--1758, doi:10.1090/S0002-9939-03-06927-2.
All eight pages were read.  Theorem~1 gives the Gibbs/BIP equivalence under
topological mixing, summable variations including finite first variation,
and finite Gurevich pressure.  In
Theorem~\ref{thm:symbolic-full-shift-pressure} these hypotheses are checked
only for the completed full shift; the positive arithmetic locus is handled
separately.

\bibitem{ConwayRadinSadun1999}
John H. Conway, Charles Radin, and Lorenzo Sadun,
\emph{On Angles Whose Squared Trigonometric Functions are Rational},
\emph{Discrete \& Computational Geometry} \textbf{22} (1999), no.~3,
321--332, doi:10.1007/PL00009463; arXiv:math-ph/9812019v1.
The versioned arXiv source was acquired to the topical shelf.  Its controlling
Plain-\TeX{} file has 42,239 bytes and SHA-256
\nolinkurl{0b37130b7e958c58c1e9f70f855f9181699b0f42d78ac5d04c29329aae0db21d}.
Theorems~1--2 are at source lines 151--168.  Lines 251--306 reconstruct the
$d=1$ St\o rmer theory from unique factorization in $\Z[i]$, orient one
Gaussian prime over each split rational prime, and prove the corresponding
angle coordinates rationally independent modulo $\pi$.  This supplies the
fixed-instance relation algorithm used in
Section~\ref{sec:chatnotes-arctangent-torsion}, not a uniform Collatz-prefix
rank theorem.

\bibitem{Calcut2009}
Jack S. Calcut,
\emph{Gaussian Integers and Arctangent Identities for $\pi$},
\emph{The American Mathematical Monthly} \textbf{116} (2009), no.~6,
515--530, doi:10.1080/00029890.2009.11920967.
The controlling 16-page PDF has 169,888 bytes and SHA-256
\nolinkurl{d650da61527c78c67caca06758363280bac7303cb2b553fc786914ddf3a4886f}.
The Main Lemma is on physical page 5 / printed page 519; Corollaries~1--3 are
on physical page 6 / printed page 520.  They prove that the relevant rational
angle is in $(\pi/4)\Z$ and that a single rational tangent gives a rational
multiple of $\pi$ only for tangent $0$ or $\pm1$.

\bibitem{GasullLucaVarona2023}
Armengol Gasull, Florian Luca, and Juan L. Varona,
\emph{Three Essays on Machin's Type Formulas},
\emph{Indagationes Mathematicae} \textbf{34} (2023), no.~6, 1373--1396,
doi:10.1016/j.indag.2023.07.002; arXiv:2302.00154v2.
The versioned arXiv source was acquired to the topical shelf.  Its controlling
source-\TeX{} file has 82,898 bytes and SHA-256
\nolinkurl{53f19ec4bf58c9e64ec51a2397401cb51fc450c7d9eb5ec61d2cf0372c162baf}.
Section~2 begins at source line 230.  Theorem~1 is at lines 246--267 and
contains $(1,1/2,1/2,3/4)$; lines 278--312 give the Gaussian root-of-unity
reformulation; and lines 341--352 derive the relevant zero double-angle
family and its exceptional $a_1=1$ specialization.  Its literal theorem scope
is two terms with nonzero right-hand side, not all zero relations or arbitrary
Collatz-prefix length.

\bibitem{Deaconu1995}
Valentin Deaconu,
\emph{Groupoids Associated with Endomorphisms},
\emph{Transactions of the American Mathematical Society} \textbf{347}
(1995), no.~5, 1779--1786.
The controlling nine-page PDF has 742,789 bytes and SHA-256
\nolinkurl{63a34e76b5404328d6eb5e7516660af357511d824a6de4b4a6286112e77cfdae}.
Physical pages 2--4 / printed pages 1779--1781 were read, and physical pages
2--3 were rendered.  They supply the set of triples, the groupoid operations,
and the compact-Hausdorff self-covering hypotheses of Theorem~1.  The
positive-odd discrete theorem in
Section~\ref{sec:chatnotes-groupoid-motivic} is proved independently because
those hypotheses do not hold there.

\bibitem{HuberWustholz2022}
Annette Huber and Gisbert W\"ustholz,
\emph{Transcendence and Linear Relations of $1$-Periods},
Cambridge Tracts in Mathematics, volume 227, Cambridge University Press,
2022.
The controlling 265-page PDF has 5,175,455 bytes and SHA-256
\nolinkurl{1743c141fd0078e7ac19805a5174078a26551d978bf6d1cb623bf3b237da3d17}.
Physical pages 90, 116--120, and 133--140 / printed pages 69, 95--99, and
112--119 were read; physical pages 119 and 138 were rendered.  Definition
8.1 fixes $1$-motives and morphisms as complexes and commutative squares,
Section~10.2 supplies the Kummer motive and path-selected logarithm, and
Proposition~12.10 supplies the $1$-period comparison.  The concrete point
$(2,3)$ and all path characters in the live chapter are independent
specializations.

\bibitem{HuberMullerStach2025}
Annette Huber and Stefan M\"uller-Stach, with contributions by Benjamin
Friedrich and Jonas von Wangenheim,
\emph{Periods and Nori Motives}, Springer, annotated version dated
1 July 2025.
The controlling 402-page PDF has 1,855,231 bytes and SHA-256
\nolinkurl{6d08014f21a616abb4ba7e177effad4d15e4a215b0f234ce3a4605c20a02b731}.
Physical pages 161--164, 231--232, and 323--324 / printed pages 139--142,
209--210, and 301--302 were read.  Definition~9.1.1 states exactly two
primitive edge types for effective pairs: contravariant functoriality from a
map of pairs and the coboundary edge from a nested triple.  It also confirms
the volume's $(k,\mathbb Q)$-Vect convention.  No path-concatenation edge is
attributed to the source.

\bibitem{BarbieriVialeKahn2016}
Luca Barbieri-Viale and Bruno Kahn,
\emph{On the Derived Category of $1$-Motives},
\emph{Ast\'erisque} \textbf{381} (2016); arXiv:1009.1900v2.
The versioned source archive was acquired to the topical shelf.  Its
controlling source file \path{der1motast-pbl.tex} has 564,248 bytes and
SHA-256
\nolinkurl{f884b90f3760ac666c1047a975e6788dfc02a1eb7f047e8eae8b57df6f013c18}.
Source lines 728--750 define a Deligne $1$-motive and its commutative-square
morphisms; lines 1529--1618 give the Cartier-dual character-lattice
machinery; lines 1736--1751 state the fully faithful curve-generated
embedding; and lines 3796--3820 define the Albanese functor and record the
integral adjunction boundary.  The source contains no Collatz-specific
$(2,3)$ specialization or Kummer character used in the live chapter.

\end{thebibliography}
\endgroup
